%%%%%%%%%%%%%%%%%%%%%%%%%%%%%%%%%%%%%%%%%%%%%%%%%%%%%%%%%%%%%%%%%%%%%%%
%%%% Таблица статуса замечаний С. П. Ковалева                     %%%%
%%%% Сборка: xelatex kovalev_status.tex                           %%%%
%%%%%%%%%%%%%%%%%%%%%%%%%%%%%%%%%%%%%%%%%%%%%%%%%%%%%%%%%%%%%%%%%%%%%%%

\documentclass[a4paper,11pt]{article}

\usepackage[a4paper,landscape,top=18mm,bottom=18mm,left=15mm,right=15mm]{geometry}

\usepackage{fontspec}
\setmainfont{Liberation Serif}
\setsansfont{Liberation Sans}
\setmonofont{Liberation Mono}
\usepackage{polyglossia}
\setdefaultlanguage{russian}
\setotherlanguage{english}

\usepackage{amsmath,amssymb}
\usepackage{longtable}
\usepackage{array}
\usepackage{tabularx}
\usepackage{xcolor}
\usepackage{booktabs}
\usepackage{microtype}
\usepackage{enumitem}

% Цвет статуса
\definecolor{statusDone}{HTML}{1A7F2F}
\newcommand{\statusDone}{{\color{statusDone}\textbf{Выполнено}}}

\usepackage{hyperref}
\hypersetup{
    unicode=true,
    pdftitle={Статус замечаний С. П. Ковалева},
    pdfauthor={К. Д. Русаков},
    colorlinks=true,
    linkcolor=blue,
    urlcolor=blue,
}

\renewcommand{\arraystretch}{1.25}
\setlength{\tabcolsep}{4pt}
\setlength{\LTcapwidth}{\textwidth}

\begin{document}

\begin{center}
{\Large\bfseries Статус замечаний С. П. Ковалева по диссертации и автореферату}\\[6pt]
{\large Работа: <<Методика распознавания и~реидентификации людей в~информационно-поисковых системах>>}\\[2pt]
Специальность 2.3.8 <<Информатика и информационные процессы>>. Автор: Русаков~К.\,Д., ИПУ~РАН.
\end{center}

\vspace{6pt}

% ============================================================
%   РАУНД 1. БАЗОВЫЕ ПРИНЦИПЫ
% ============================================================
\section*{Раунд 1. Базовые принципы}

\begin{longtable}{@{}p{0.025\textwidth}p{0.33\textwidth}p{0.13\textwidth}p{0.46\textwidth}@{}}
\toprule
\textbf{№} & \textbf{Замечание} & \textbf{Статус} & \textbf{Где / комментарий}\\
\midrule
\endfirsthead
\toprule
\textbf{№} & \textbf{Замечание} & \textbf{Статус} & \textbf{Где / комментарий}\\
\midrule
\endhead
\bottomrule
\endfoot

1 & \textbf{Глобальная абстракция.} Перевести постановку с <<реидентификации людей>> на <<распознавание сложных многокомпонентных структур по их разрозненным частям>>.
& \statusDone
& Введение диссертации и общая характеристика автореферата переписаны через рамку <<многокомпонентная структура $C_1,\ldots,C_K$>>; реидентификация представлена как наиболее представительный частный случай.\\
\midrule

2 & \textbf{Сначала формулы, потом люди.} Сперва изложить универсальный методический аппарат (детектор, кодировщики ViT с~ArcFace, байесовская модель с~MAP-критерием) в математическом вакууме, потом конкретизировать на реидентификацию людей.
& \statusDone
& Во введении диссертации сформирован блок <<Формальная постановка задачи>> для произвольного числа компонентов~$K$. В главе~2 (подраздел~2.2.3) сначала идет параграф <<Общий случай: произвольное число компонентов структуры>>, затем конкретизация на~$K=2$.\\
\midrule

3 & \textbf{Доказательство универсальности.} Подсветить опубликованные работы по смежным предметным областям (коррозия, номера, БПЛА) как подтверждение применимости метода вне реидентификации людей.
& \statusDone
& В главе~3 добавлен раздел <<Масштабируемость метода на смежные предметные области>> со ссылками на опубликованные работы автора по распознаванию дефектов металлоконструкций, идентификации транспортных средств, мониторингу с~БПЛА, антиспуфингу и прогнозу технического состояния космических аппаратов.\\
\end{longtable}

\vspace{4pt}

% ============================================================
%   РАУНД 1. ТЕХНИЧЕСКИЕ УТОЧНЕНИЯ
% ============================================================
\section*{Раунд 1. Технические уточнения}

\begin{longtable}{@{}p{0.025\textwidth}p{0.33\textwidth}p{0.13\textwidth}p{0.46\textwidth}@{}}
\toprule
\textbf{№} & \textbf{Замечание} & \textbf{Статус} & \textbf{Где / комментарий}\\
\midrule
\endfirsthead
\toprule
\textbf{№} & \textbf{Замечание} & \textbf{Статус} & \textbf{Где / комментарий}\\
\midrule
\endhead
\bottomrule
\endfoot

a & Привести байесовскую интеграцию в соответствие с программной реализацией: убрать избыточную формулу с коэффициентом~$\alpha$, описать механизм автоматического взвешивания через ковариационные матрицы~$\Sigma_y$ (свидетельство о регистрации программы для ЭВМ № 2026617996 от 23.03.2026).
& \statusDone
& В автореферате (после уравнения MAP-критерия) и в главе~2 диссертации дополнительный гиперпараметр комбинирования удален; явно показано, что роль коэффициентов доверия играют именно ковариационные матрицы $\Sigma_{y,k}$ гауссовых правдоподобий.\\
\midrule

b & Добавить в третью главу раздел <<Масштабируемость метода>> со ссылками на работы по коррозии, номерам, БПЛА.
& \statusDone
& В главе~3 диссертации добавлен соответствующий раздел; пересекается с пунктом~3 раунда~1.\\
\midrule

c & Добавить характеристики серверной инфраструктуры экспериментов.
& \statusDone
& В главе~3 и в автореферате приведена таблица характеристик: Ubuntu 22.04, AMD Ryzen 9 7950X, два графических процессора NVIDIA RTX 3090 Ti, 128 ГБ оперативной памяти, CUDA 12.5.\\
\midrule

d & Восстановить определения переменных в формулах.
& \statusDone
& Каждая основная формула во введении, в главе~2 и в автореферате сопровождается списком определений всех переменных: $\mathcal{L}_0$, $N_{\text{pred}}$, $\mathbb{I}_{\text{inner}}$, $\mathcal{L}_{\text{inner}}$, $\lambda_{\text{inner}}$, $\boldsymbol{\mu}_{y,k}$, $\Sigma_{y,k}$ и др.\\
\midrule

e & Обосновать выбор архитектуры YOLOv8 (блоки C2f, SPPF).
& \statusDone
& В автореферате и в главе~2 диссертации явно указано на <<сбалансированное сочетание высокой скорости обработки и точности локализации за счет эффективных блоков C2f и механизма Spatial Pyramid Pooling Fast (SPPF)>>.\\
\midrule

f & Упомянуть Акт о внедрении в ФГБУ <<НИЦ <<Институт имени Н.\,Е.\,Жуковского>>>> (сокращение времени принятия решения до~5\%).
& \statusDone
& Внесено в общую характеристику автореферата (апробация и практическая значимость), а также во введение и заключение диссертации. Указана дата Акта о внедрении (10.06.2024), приведены ссылки на свидетельства о государственной регистрации программ для ЭВМ.\\
\midrule

g & Исправить оформительские пометки в списке публикаций.
& \statusDone
& В списке публикаций автореферата убраны незаполненные заглушки диапазона страниц в двух публикациях (конференция MLSD, 2024 и 2023); написание <<Юго-Западного Государственного университета>> исправлено на <<Юго-Западного государственного университета>> (две позиции).\\
\midrule

h & Унифицировать терминологию (<<многокомпонентная структура>>), оформление по ГОСТ.
& \statusDone
& Терминология <<многокомпонентная структура>> используется сквозь все разделы введения, обеих частей автореферата и глав диссертации. Удалены рудиментарные обороты <<топологически связанных частей объекта>>.\\
\end{longtable}

\vspace{4pt}

% ============================================================
%   РАУНД 2
% ============================================================
\section*{Раунд 2. Замечания после переработки введения}

\begin{longtable}{@{}p{0.025\textwidth}p{0.33\textwidth}p{0.13\textwidth}p{0.46\textwidth}@{}}
\toprule
\textbf{№} & \textbf{Замечание} & \textbf{Статус} & \textbf{Где / комментарий}\\
\midrule
\endfirsthead
\toprule
\textbf{№} & \textbf{Замечание} & \textbf{Статус} & \textbf{Где / комментарий}\\
\midrule
\endhead
\bottomrule
\endfoot

1 & Тезис о масштабируемости метода на другие предметные области не подтвержден в основном тексте.
& \statusDone
& В главе~3 добавлен раздел <<Масштабируемость метода на смежные предметные области>>, с подразделами по работам автора (коррозия, номера, мониторинг с~БПЛА, антиспуфинг, космические аппараты).\\
\midrule

2 & Индексы <<face>>/<<лицо>> в формуле потерь YOLOv8 специфичны для распознавания лиц; для соответствия абстрактной постановке требуется обобщить ($\mathcal{L}_{\text{inner}}$, $\lambda_{\text{inner}}$ — <<вложенный компонент>>), интерпретацию <<вложенный компонент~$=$~лицо>> дать отдельно как частный случай.
& \statusDone
& Во введении, в главе~2 и в автореферате обозначения обобщены: $\mathcal{L}_{\text{face}}$, $\lambda_{\text{face}}$, $\mathbb{I}_{\text{face}}$ заменены на $\mathcal{L}_{\text{inner}}$, $\lambda_{\text{inner}}$, $\mathbb{I}_{\text{inner}}$ (<<вложенный компонент>>); соответствующие переменные переименованы и в программной реализации функции потерь. Конкретные обозначения для лица ($\Phi_{\text{face}}$, $\mathrm{ROI}_{\text{face}}$) оставлены только там, где речь именно о лицевой модальности как одной из двух частей в~случае~$K=2$.\\
\midrule

3 & Все схемы процессов и алгоритмов содержат специфику реидентификации людей; нужен обобщенный вариант для произвольного $K$.
& \statusDone
& Подготовлены обобщенные блок-схемы: схема информационных процессов информационно-поисковой системы и блок-схема байесовской интеграции для произвольного~$K$. Они включены рядом с прежними рисунками в автореферате и в диссертации; подписи существующих рисунков переведены в формат <<Частный случай ($K=2$): силуэт-лицо>>.\\
\midrule

4 & Байесовская модель в гл.~2 описана для пары признаков лица и силуэта; абстрактный многокомпонентный случай произвольного набора доступных компонентов $\mathcal{M}(t)\subseteq\{1,\ldots,K\}$ во введении дан, но в основном тексте не раскрыт.
& \statusDone
& В подраздел~2.2.3 <<Байесовская модель интеграции>> главы~2 добавлен параграф <<Общий случай: произвольное число компонентов структуры>> с соответствующими уравнениями (правило Байеса, MAP-критерий, гауссова модель правдоподобий) для произвольного~$K$; ниже идет конкретизация на~$K=2$.\\
\midrule

5 & Пункт 3 научной новизны об абстрактном методе интеграции не раскрыт в гл.~2.
& \statusDone
& Тот же параграф <<Общий случай>> в подразделе~2.2.3 раскрывает абстрактный метод интеграции для произвольного~$K$ и явно соединяет его с положением о байесовской модели интеграции (положение~2), выносимым на защиту.\\
\midrule

6 & Название работы <<Комплексная методика реидентификации людей в~информационно-поисковых системах>> расходится с продекларированной абстрактной постановкой; рекомендована переформулировка (<<Методика распознавания многокомпонентных структур в~информационно-поисковых системах (на примере реидентификации людей)>>).
& \statusDone
& \textbf{Название приведено к~формулировке Акта о~внедрении} в~ФГБУ <<НИЦ <<Институт имени Н.\,Е.\,Жуковского>>>>: <<Методика распознавания и~реидентификации людей в~информационно-поисковых системах>> (ранее <<Комплексная методика реидентификации людей\ldots>>). Полная переформулировка к~абстрактной форме <<распознавания многокомпонентных структур>> не~выполнена, так как название фиксируется Актом о~внедрении. Расхождение между названием и абстрактной рамкой компенсируется явной формулировкой во введении и в~автореферате о том, что реидентификация людей выступает как наиболее представительный частный случай предлагаемой методики, а в~главе~2 (общий случай~$K$) и~главе~3 (раздел масштабируемости) развернут универсальный аппарат.\\
\midrule

7 & Термин <<множество идентичностей>> режет ухо; заменить на <<множество классов идентификации>>~/ <<множество личностей>> по контексту.
& \statusDone
& Заменено во всех разделах автореферата и диссертации (около 16~вхождений). Используется <<множество классов идентификации>>~/ <<личностей>>~/ <<классов>> по контексту.\\
\end{longtable}

\vspace{4pt}

% ============================================================
%   РАУНД 3
% ============================================================
\section*{Раунд 3. Замечания от 5 июня 2026 г.}

\begin{longtable}{@{}p{0.025\textwidth}p{0.33\textwidth}p{0.13\textwidth}p{0.46\textwidth}@{}}
\toprule
\textbf{№} & \textbf{Замечание} & \textbf{Статус} & \textbf{Где / комментарий}\\
\midrule
\endfirsthead
\toprule
\textbf{№} & \textbf{Замечание} & \textbf{Статус} & \textbf{Где / комментарий}\\
\midrule
\endhead
\bottomrule
\endfoot

1 & Абстракция выполнена фрагментарно; п.~2 научной новизны, положений на защиту и заключения переформулировать; основной текст переработать (всюду, кроме разделов 1.2, 2.2.3 и 3.4, речь только о людях).
& \statusDone\newline {\footnotesize (по гл.~1 в основном)}
& П.~2 и остальные пункты научной новизны, положений на защиту и заключения переписаны от абстрактного к частному: ведущая сущность~-- распознавание многокомпонентных структур по их разрозненным частям, реидентификация людей (силуэт-лицо, $K=2$)~-- наиболее представительный частный случай и апробация. Основной текст глав~2 и~3 переведен на сквозную рамку многокомпонентных структур; конкретика реидентификации вынесена как частный случай. В главе~1 рамка задана в постановке задачи для произвольного~$K$, в обобщенной схеме процессов и в выводах; обзор конкретных методов распознавания лиц (раздел~1.1: метод главных компонент, линейный дискриминантный анализ, метод Виолы-Джонса, гистограммы ориентированных градиентов) предметно-специфичен по своей природе.\\
\midrule

2 & В разделе~3.4 не ограничиваться общими рассуждениями~-- привести численные значения показателей качества распознавания многокомпонентных структур в нескольких смежных задачах.
& \statusDone
& В раздел~3.4 добавлены численные показатели для двух смежных задач, реально устроенных покомпонентно. (1)~Распознавание автомобильного регистрационного знака как структуры из~символов~-- по опубликованной работе автора (модульная система распознавания номеров, конференция MLSD, 2020): детектирование знака mAP[0,5:0,95]~81,2\,\%, детектирование символов~95,2\,\%, распознавание символов~99,32\,\%. (2)~Мультикамерная оценка движения БПЛА как интеграция частичных наблюдений~-- по работе автора (оценка скорости и ускорения БПЛА мультикамерным методом, MLSD, 2023); приведенное число характеризует точность оценивания кинематики, а не классификационное качество. Коррозия, антиспуфинг и прогноз состояния космического аппарата оставлены как качественные свидетельства широты применимости, без приписывания им предложенного метода; отдельный эксперимент специально под смежную задачу не~ставился.\\
\midrule

3 & В сравнительной таблице (Таблица~4 автореферата) неясно, почему на разных наборах данных использовались разные методы для сравнения с предложенным.
& \statusDone
& После сравнительной таблицы методов (в главе~3 диссертации и в автореферате) добавлено пояснение: большинство методов отчитывается лишь на одном-двух наборах из CUHK03, Market-1501, MARS (прочерк означает отсутствие опубликованного значения, а не низкое качество); наборы относятся к покадровым (CUHK03, Market-1501) и видео (MARS); значения для CUHK03 получены при разных протоколах оценки (ранний протокол одиночного предъявления против современного протокола с раздельной отчетностью по ручной разметке и по детектированным рамкам); предложенный метод оценен в закрытой постановке распознавания, а большинство сравниваемых~-- в открытом ранжирующем поиске. Дополнительно введены опорные методы, отчитывающиеся одновременно на CUHK03 и Market-1501 (MGN, BDB, MHN, Auto-ReID, Pyramid, RGA-SC, SCSN), и метод AGW, отчитывающийся сразу на всех трех наборах.\\
\midrule

4 & В диссертации привести ссылки на литературу по каждому из методов первой колонки сравнительной таблицы.
& \statusDone
& К каждому методу первой колонки таблицы проставлена ссылка на литературу; для этого в список литературы добавлено 28 верифицированных источников (20 методов сравнения по наборам, один базовый комбинированный метод и 7 опорных методов для межнаборного сопоставления); еще один источник (метод AGW, отчитывающийся на всех трех наборах) уже присутствовал. Все источники проверены по первоисточникам (arXiv, материалы конференций CVF, AAAI, IEEE, ACM, журнал IEEE TPAMI). Перекрестные ссылки в диссертации разрешаются без ошибок.\\
\end{longtable}

\end{document}
